\documentclass[11pt,oneside]{article}

\usepackage[T1]{fontenc}
\usepackage[utf8]{inputenc}
\usepackage{lmodern}
\usepackage{microtype}
\usepackage{geometry}
\geometry{margin=1in,headheight=15pt}
\usepackage{amsmath,amssymb,amsthm,mathtools}
\usepackage{booktabs,longtable,tabularx,array}
\usepackage{enumitem}
\usepackage{xcolor}
\usepackage{xurl}
\usepackage{hyperref}
\usepackage{listings}
\usepackage{fancyhdr}
\usepackage{titlesec}
\usepackage{needspace}

\definecolor{theoremblue}{RGB}{29,78,121}
\definecolor{newgreen}{RGB}{18,103,75}
\definecolor{warningred}{RGB}{148,46,42}
\definecolor{computeviolet}{RGB}{92,61,145}
\definecolor{softgray}{RGB}{245,246,247}
\definecolor{rulegray}{RGB}{96,104,112}

\hypersetup{
  colorlinks=true,
  linkcolor=theoremblue,
  citecolor=newgreen,
  urlcolor=theoremblue,
  pdftitle={Alaoglu--Erdos Progress: Mesoscopic Factorial Cocycles and Prime-Window Obstructions},
  pdfauthor={OpenAI GPT-5.6 Pro, prepared for Vladimir Reshetnikov},
  pdfsubject={A new factorial-cocycle attack on the Alaoglu--Erdos integer-exponent conjecture}
}
\urlstyle{same}

\pagestyle{fancy}
\fancyhf{}
\lhead{Alaoglu--Erd\H{o}s: factorial-cocycle supplement}
\rhead{August 22, 2026}
\cfoot{\thepage}
\renewcommand{\headrulewidth}{0.4pt}

\titleformat{\section}{\Large\bfseries}{\thesection}{0.75em}{}
\titleformat{\subsection}{\large\bfseries}{\thesubsection}{0.75em}{}
\titleformat{\subsubsection}{\normalsize\bfseries}{\thesubsubsection}{0.75em}{}

\setlist[itemize]{leftmargin=2em,itemsep=0.3em,topsep=0.45em}
\setlist[enumerate]{leftmargin=2.3em,itemsep=0.3em,topsep=0.45em}
\setlength{\parindent}{0pt}
\setlength{\parskip}{0.55em}
\setlength{\emergencystretch}{3em}

\newtheorem{theorem}{Theorem}[section]
\newtheorem{proposition}[theorem]{Proposition}
\newtheorem{lemma}[theorem]{Lemma}
\newtheorem{corollary}[theorem]{Corollary}
\newtheorem{conjecture}[theorem]{Conjecture}
\newtheorem{criterion}[theorem]{Criterion}
\newtheorem{target}[theorem]{Research target}
\theoremstyle{definition}
\newtheorem{definition}[theorem]{Definition}
\newtheorem{experiment}[theorem]{Exact experiment}
\newtheorem{warning}[theorem]{Audit warning}
\theoremstyle{remark}
\newtheorem{remark}[theorem]{Remark}
\newtheorem{observation}[theorem]{Observation}

\newcommand{\N}{\mathbb N}
\newcommand{\Nzero}{\mathbb N_0}
\newcommand{\Z}{\mathbb Z}
\newcommand{\Q}{\mathbb Q}
\newcommand{\R}{\mathbb R}
\newcommand{\thetaAE}{\vartheta}
\newcommand{\den}{\operatorname{den}}
\newcommand{\ord}{\operatorname{ord}}
\newcommand{\rad}{\operatorname{rad}}
\newcommand{\vp}{v_p}
\newcommand{\fall}[2]{\left(#1\right)_{\!\underline{#2}}}
\newcommand{\EF}{\mathsf E}
\newcommand{\cF}{\mathcal F}
\newcommand{\cB}{\mathcal B}
\newcommand{\cP}{\mathcal P}
\newcommand{\statusnew}{\textcolor{newgreen}{\textbf{[proved in this report]}}}
\newcommand{\statuscomp}{\textcolor{computeviolet}{\textbf{[exact computation]}}}
\newcommand{\statusext}{\textcolor{theoremblue}{\textbf{[external theorem]}}}
\newcommand{\statusopen}{\textcolor{warningred}{\textbf{[open bridge]}}}

\newenvironment{statusbox}[1]
{\begin{center}\begin{minipage}{0.94\textwidth}\raggedright\hrule\vspace{0.45em}
\textbf{Status.} #1\par\vspace{0.35em}}
{\vspace{0.35em}\hrule\end{minipage}\end{center}}

\newenvironment{resultbox}
{\begin{center}\begin{minipage}{0.94\textwidth}\hrule\vspace{0.45em}}
{\vspace{0.35em}\hrule\end{minipage}\end{center}}

\lstdefinestyle{pythonstyle}{
  language=Python,
  basicstyle=\ttfamily\scriptsize,
  backgroundcolor=\color{softgray},
  frame=single,
  rulecolor=\color{rulegray},
  numbers=left,
  numberstyle=\tiny\color{rulegray},
  numbersep=8pt,
  showstringspaces=false,
  breaklines=true,
  columns=fullflexible,
  keepspaces=true,
  tabsize=4
}

\begin{document}
\hypersetup{pageanchor=false}
\pagenumbering{gobble}

\begin{titlepage}
\centering
\vspace*{0.65cm}
{\Huge\bfseries The Alaoglu--Erd\H{o}s Conjecture\par}
\vspace{0.28cm}
{\LARGE\bfseries Mesoscopic Factorial Cocycles,\par}
\vspace{0.12cm}
{\LARGE\bfseries Prime-Window Obstructions, and a Linear-Width Escape Target\par}
\vspace{0.85cm}
{\large A nonduplicative supplemental progress report\par}
\vspace{0.1cm}
{\large to the unified ProveIt report supplied on August 22, 2026\par}
\vfill
\begin{minipage}{0.91\textwidth}
\small
\textbf{Bottom line.}
I did not obtain a complete proof of
\[
 2^x,3^x\in\Z\quad\Longrightarrow\quad x\in\Z.
\]
I did obtain a new, exact attack that is not developed in the supplied report, together with
an essentially complete diagnosis of its fixed-width and subcritical-width forms.
For a hypothetical counterexample put $M=2^x$, $A=3^x$, and
$N_n=A^n$, $K_n=M^n$.  The universal identity
\[
 K_n=N_n^{\log 2/\log 3},\qquad
 \frac12<\frac{\log2}{\log3}<\frac23
\]
places the falling-factorial probability
\[
 R_n=\frac{(A^n)_{M^n}}{(A^n)^{M^n}}
\]
in a uniquely favorable ``birthday'' regime: pair collisions diverge as
$(M^2/A)^n=(4/3)^{xn}$, while triple collisions vanish as
$(M^3/A^2)^n=(8/9)^{xn}$.  A finite shift polynomial can cancel the sole growing
collision mode exactly.

For every integer polynomial $P$, the signed factorial cocycle
\[
 \cF_{P,n}=\prod_j\left(\frac{A^{n+j}!}{(A^{n+j}-M^{n+j})!}\right)^{[T^j]P}
\]
tends to $1$ precisely when
\[
 P(M)=P'(M)=P(M^2/A)=0.
\]
The primitive minimal annihilator is
\[
 P_*(T)=(T-M)^2\left(\frac{A}{g}T-\frac{M^2}{g}\right),
 \qquad g=\gcd(A,M^2),
\]
and its cocycle satisfies the sharp asymptotic
\[
 \log\cF_{P_*,n}\sim
 \frac{M^2(A-M)}{6gA}
 \left(M-\frac{M^3}{A^2}\right)^2
 \left(\frac{M^3}{A^2}\right)^n>0.
\]
Thus it is a positive rational number converging exponentially to $1$.

The arithmetic side is hostile in an equally exact way.  Using the inherited exclusion
$x\ge12$, Huxley's prime number theorem in short intervals, and Brun--Titchmarsh covering,
I prove that every such cocycle whose shifts extend less than
\[
 \bigl(\log_2 3-1-\varepsilon\bigr)n
 =\bigl(0.5849625007\ldots-\varepsilon\bigr)n
\]
beyond its highest negative coefficient has a reduced denominator with logarithm
$\gg M^n$.  In particular, every fixed-shift factorial or binomial cocycle fails by an
exponentially large prime-window denominator.  The first possible escape is therefore not
``improve a constant'' but a sharply specified linear-width prime-transport problem.  I give
its exact valuation inequalities, two phase thresholds, a finite optimization formulation,
and a Lean-oriented decomposition.
\end{minipage}
\vfill
{\large Prepared for Vladimir Reshetnikov\par}
\vspace{0.12cm}
{\large August 22, 2026\par}
\end{titlepage}
\hypersetup{pageanchor=true}

\pagenumbering{roman}

\begin{abstract}
Let $x>0$ and suppose that $M=2^x$ and $A=3^x$ are integers.  The supplied unified report
already develops the exact conditional solution monoid, local and global transcendence
barriers, generalized Vandermonde determinants, $p$-adic and Kummer diagnostics, compact-orbit
models, entire interpolation, Stieltjes moments, carry languages, and numerous quantified
no-go theorems.  This supplement deliberately avoids repeating those constructions.

The new object is the mesoscopic falling factorial
$F_n=(A^n)_{M^n}$ and its normalized no-collision probability
$R_n=F_n/(A^n)^{M^n}$.  Since
$\alpha=\log M/\log A=\log2/\log3\in(1/2,2/3)$, one has
$M^{2n}/A^n\to\infty$ but $M^{3n}/A^{2n}\to0$.  An explicit logarithmic expansion gives
\[
 \log F_n=nM^n\log A-\frac12\left(\frac{M^2}{A}\right)^n
 +\frac12\left(\frac MA\right)^n
 -\frac16\left(\frac{M^3}{A^2}\right)^n
 +O\!\left(\left(\frac{M^4}{A^3}\right)^n\right).
\]
For a polynomial $P(T)=\sum c_jT^j$, the cocycle
$\cF_{P,n}=\prod_jF_{n+j}^{c_j}$ therefore tends to one if and only if
$P$ has a double zero at $M$ and a zero at $M^2/A$.  This yields a canonical minimal
near-unit and a complete fixed-shift classification.

The denominator is analyzed prime by prime through exact factorial-valuation floor sums.
Every annihilator has at least three coefficient sign variations.  If $k$ is its highest
negative shift, primes in the interval
$(A^{n+k}-M^{n+k},A^{n+k}]$ enter the denominator unless a higher positive factorial interval
captures a multiple.  A Brun--Titchmarsh covering estimate shows that, for the least
hypothetical counterexample (hence $x\ge12$), higher shifts of total span below
$(\log_2 3-1-\varepsilon)(n+k)$ can capture less than a fixed fraction of the short-interval
prime mass.  Huxley's theorem supplies total logarithmic prime mass asymptotic to $M^{n+k}$,
so the reduced denominator has exponential logarithmic size.  The same obstruction applies
to binomial cocycles; their near-unit classification adds a double zero at $1$.

The resulting research target is concrete: construct a linearly widening signed factorial
ratio, satisfying three archimedean annihilation constraints and all prime-power valuation
inequalities, with shift span at least $(\log_2 3-1)n$ but sufficiently small coefficient
height.  This is expressed as an exact integer linear feasibility problem and separated into
finite, analytic, and formalizable layers.
\end{abstract}

\tableofcontents
\clearpage
\pagenumbering{arabic}

\section{Scope, trust boundary, and nonduplication}
\label{sec:scope}

\subsection{What is inherited}

The attached unified report is treated as the baseline, not as material to be rewritten\cite{UnifiedReport2026}.
In particular, this supplement imports the following established facts and conventions:
\begin{enumerate}
\item a putative nonintegral solution is positive and transcendental;
\item after passing to the least nonintegral solution one has an exact primitive output pair
      $(M,A)=(2^x,3^x)$ and an exact rank-two conditional solution monoid;
\item the verified finite exclusion $2^x<4096$ implies that a least nonintegral solution
      satisfies $x\ge12$;
\item the two-base problem is an exact $2\times2$ four-exponentials boundary, so no assertion
      below silently invokes four exponentials or Schanuel;
\item determinant, $p$-adic, Kummer, compact-orbit, rational-function, entire-interpolation,
      moment, and carry-language routes already have detailed audits in the baseline report.
\end{enumerate}

Only item 3 enters the strongest new denominator theorem.  The elementary factorial
classification itself holds for every positive integral pair $1<M<A$ and does not use the
real-power origin of that pair.  The historical conjecture itself originates in the work of
Alaoglu and Erd\H{o}s on highly composite and related numbers\cite{AlaogluErdos1944}.

\subsection{What is new here}

The baseline report does discuss factorial divisors of interpolation determinants and Smith
profiles.  Those are not the construction below.  Here factorials arise from the occupancy
product
\[
 (N)_K=N(N-1)\cdots(N-K+1),
\]
with the special geometric scaling $(N,K)=(A^n,M^n)$.  The central new features are:
\begin{enumerate}
\item the exact mesoscopic window $K=N^{\alpha}$ with $1/2<\alpha<2/3$;
\item a shift-polynomial calculus that classifies every fixed factorial near-unit;
\item the primitive minimal annihilator $P_*$ and its exact leading defect;
\item a prime-window theorem forcing denominator mass for all fixed and subcritical
      linearly widening shifts;
\item the universal escape width $\log_2 3-1$ and a finite valuation-feasibility target.
\end{enumerate}

\begin{statusbox}{\statusnew{} for all statements proved in Sections
\ref{sec:mesoscopic}--\ref{sec:primewindows}, except for the explicitly cited short-interval
prime theorem and Brun--Titchmarsh inequality.  Numerical tables are \statuscomp{}.}
\end{statusbox}

\subsection{Competitive-status note}

The public material located during this run did not contain an inspectable proof text or Lean
repository for the Alaoglu--Erd\H{o}s conjecture.  The reported competing route is therefore
treated as unavailable information: no mathematical step below is inferred from it, and every
claim in this report is auditable independently of the rumor.

\section{The mesoscopic exponent forced by the bases two and three}
\label{sec:mesoscopic}

Assume throughout that
\begin{equation}
 M=2^x\in\N,\qquad A=3^x\in\N,
 \qquad x>0.
 \label{eq:MA}
\end{equation}
Then $1<M<A$ and
\begin{equation}
 \alpha:=\frac{\log M}{\log A}=\frac{\log2}{\log3},
 \qquad
 \thetaAE:=\alpha^{-1}=\log_2 3.
 \label{eq:alpha-theta}
\end{equation}
The two elementary inequalities $4>3$ and $8<9$ give
\begin{equation}
 \frac12<\alpha<\frac23,
 \qquad
 1<\thetaAE<2.
 \label{eq:alpha-window}
\end{equation}
Numerically,
\begin{align}
 \alpha&=0.6309297535714574\ldots,\notag\\
 1-\alpha&=0.3690702464285426\ldots,\notag\\
 \thetaAE-1&=0.5849625007211563\ldots,\notag\\
 2-\thetaAE&=0.4150374992788437\ldots.
 \label{eq:universal-constants}
\end{align}

Put
\begin{equation}
 N_n=A^n,
 \qquad K_n=M^n=N_n^{\alpha},
 \qquad q:=\frac MA=\left(\frac23\right)^x.
 \label{eq:NKq}
\end{equation}
Three ratios control everything that follows:
\begin{equation}
 \lambda:=\frac{M^2}{A}=\left(\frac43\right)^x>1,
 \qquad
 \eta:=\frac{M^3}{A^2}=\left(\frac89\right)^x<1,
 \qquad
 \tau:=\frac{M^4}{A^3}=\left(\frac{16}{27}\right)^x<q<\eta.
 \label{eq:lambda-eta-tau}
\end{equation}
Thus
\begin{equation}
 \frac{K_n^2}{N_n}=\lambda^n\longrightarrow\infty,
 \qquad
 \frac{K_n^3}{N_n^2}=\eta^n\longrightarrow0.
 \label{eq:mesoscopic-two-three}
\end{equation}

\begin{observation}[The one-divergent-cumulant regime]
\label{obs:one-divergent}
For occupancy of $K_n$ labeled balls in $N_n$ boxes, the pair-collision scale diverges while
the triple-collision scale vanishes.  This happens because the fixed exponent
$\alpha=\log2/\log3$ lies strictly between $1/2$ and $2/3$.  No estimate involving the
unknown size of $x$ is needed.
\end{observation}

The inherited finite exclusion becomes useful through $q$.  For a least nonintegral
solution, $x\ge12$, and therefore
\begin{equation}
 q\le q_{12}:=\left(\frac23\right)^{12}
 =0.007707346629258934\ldots.
 \label{eq:q12}
\end{equation}
This small ratio will leave almost all short-interval primes uncaptured by higher factorial
windows.

\section{The birthday product and its exact asymptotic expansion}
\label{sec:birthday}

\subsection{Falling factorial and no-collision probability}

Define
\begin{equation}
 F_n:=\fall{A^n}{M^n}
 =\frac{A^n!}{(A^n-M^n)!}\in\N,
 \qquad
 R_n:=\frac{F_n}{(A^n)^{M^n}}.
 \label{eq:Fn-Rn}
\end{equation}
Equivalently,
\begin{equation}
 R_n=\prod_{j=0}^{M^n-1}\left(1-\frac{j}{A^n}\right).
 \label{eq:R-product}
\end{equation}
It is the probability that $M^n$ independent uniform samples from an $A^n$-element set are
all distinct.  The probability language is only mnemonic; all arguments use the exact
finite product.

\subsection{Two-term expansion with a geometric remainder}

\begin{lemma}[Uniform logarithmic remainder]
\label{lem:log-remainder}
For $0\le u<1$,
\begin{equation}
 0\le -\log(1-u)-u-\frac{u^2}{2}
 \le \frac{u^3}{3(1-u)}.
 \label{eq:log-remainder}
\end{equation}
\end{lemma}

\begin{proof}
Expand $-\log(1-u)=\sum_{r\ge1}u^r/r$.  The remaining terms are nonnegative, and for
$r\ge3$ one has $1/r\le1/3$, giving the geometric upper bound.
\end{proof}

\begin{proposition}[Exact mesoscopic expansion]
\label{prop:Fn-expansion}
As $n\to\infty$,
\begin{align}
 \log F_n
 &=nM^n\log A
   -\frac12\lambda^n
   +\frac12q^n
   -\frac16\eta^n
   +O(\tau^n).
 \label{eq:Fn-expansion}
\end{align}
More explicitly,
\begin{align}
 \log F_n
 &=nM^n\log A
 -\frac{M^n(M^n-1)}{2A^n}
 -\frac{M^n(M^n-1)(2M^n-1)}{12A^{2n}}
 +E_n,
 \label{eq:Fn-expansion-exact}
\end{align}
where, once $q^n\le1/2$,
\begin{equation}
 |E_n|\le
 \frac{M^{4n}}{12A^{3n}(1-q^n)}.
 \label{eq:En-bound}
\end{equation}
\end{proposition}

\begin{proof}
Take logarithms in \eqref{eq:R-product} and apply Lemma~\ref{lem:log-remainder} with
$u=j/A^n$.  The first two power sums are
\[
 \sum_{j=0}^{K-1}j=\frac{K(K-1)}2,
 \qquad
 \sum_{j=0}^{K-1}j^2=\frac{K(K-1)(2K-1)}6.
\]
For the remainder, $j/A^n\le q^n$ and
$\sum_{j<K}j^3=[K(K-1)/2]^2<K^4/4$, which gives
\eqref{eq:En-bound}.  Expanding the two displayed polynomials in $K=M^n$ gives
\begin{align*}
 -\frac{K(K-1)}{2N}
 &=-\frac12\lambda^n+\frac12q^n,\\
 -\frac{K(K-1)(2K-1)}{12N^2}
 &=-\frac16\eta^n+\frac14q^{2n}
   -\frac1{12}\left(\frac{M}{A^2}\right)^n.
\end{align*}
Both omitted bases are smaller than $\tau=M^4/A^3$, while
\eqref{eq:En-bound} is $O(\tau^n)$.
\end{proof}

\begin{remark}
The expansion is unusually clean: $\lambda^n$ is the only growing correction to the bulk
term $nM^n\log A$.  Every collision mode of order at least three decays.  This is the exact
asymmetry that the shift operator will exploit.
\end{remark}

\section{Shift-polynomial calculus and complete fixed-shift classification}
\label{sec:shift}

\subsection{Polynomial shifts}

Let $\EF$ denote the forward shift on sequences,
$(\EF u)_n=u_{n+1}$.  For
\begin{equation}
 P(T)=\sum_{j=0}^{d}c_jT^j\in\Z[T]
 \label{eq:P-def}
\end{equation}
define
\begin{equation}
 (P(\EF)u)_n:=\sum_{j=0}^{d}c_ju_{n+j}
 \label{eq:shift-eval}
\end{equation}
and the positive rational factorial cocycle
\begin{equation}
 \cF_{P,n}:=\prod_{j=0}^{d}F_{n+j}^{c_j}.
 \label{eq:factorial-cocycle}
\end{equation}
Negative coefficients mean reciprocals.  The logarithm linearizes the construction:
\begin{equation}
 \log\cF_{P,n}=P(\EF)(\log F)_n.
 \label{eq:cocycle-log}
\end{equation}

Two elementary identities drive the classification.

\begin{lemma}[Shift evaluation on exponential-polynomial modes]
\label{lem:shift-modes}
For every $\rho>0$,
\begin{align}
 P(\EF)(\rho^n)&=P(\rho)\rho^n,
 \label{eq:shift-exp}\\
 P(\EF)(n\rho^n)&=\rho^n\bigl(nP(\rho)+\rho P'(\rho)\bigr).
 \label{eq:shift-nexp}
\end{align}
\end{lemma}

\begin{proof}
The first identity is immediate.  For the second,
\[
 \sum_jc_j(n+j)\rho^{n+j}
 =\rho^n\left(n\sum_jc_j\rho^j+
                  \sum_jjc_j\rho^j\right)
 =\rho^n\bigl(nP(\rho)+\rho P'(\rho)\bigr).
\]
\end{proof}

\begin{theorem}[Master asymptotic for a fixed shift polynomial]
\label{thm:master-asymptotic}
For fixed $P\in\Z[T]$,
\begin{align}
 \log\cF_{P,n}
 &=\log A\,M^n\bigl(nP(M)+MP'(M)\bigr)
   -\frac12P(\lambda)\lambda^n
   +\frac12P(q)q^n\notag\\
 &\hspace{3em}
   -\frac16P(\eta)\eta^n
   +O_P(\tau^n).
 \label{eq:master-asymptotic}
\end{align}
\end{theorem}

\begin{proof}
Apply Lemma~\ref{lem:shift-modes} term by term to
Proposition~\ref{prop:Fn-expansion}.  Since $P$ is fixed, shifting a geometric
$O(\tau^n)$ remainder through the finitely many indices $0,\ldots,d$ preserves its order.
\end{proof}

\subsection{Every fixed factorial near-unit}

\begin{theorem}[Near-unit classification]
\label{thm:near-unit-classification}
Let $P\in\Z[T]$.  Then
\begin{equation}
 \cF_{P,n}\longrightarrow1
 \qquad(n\to\infty)
 \label{eq:near-unit-limit}
\end{equation}
if and only if
\begin{equation}
 P(M)=0,
 \qquad P'(M)=0,
 \qquad P\!\left(\frac{M^2}{A}\right)=0.
 \label{eq:three-annihilation-conditions}
\end{equation}
Equivalently, $M$ is a double root of $P$ and $M^2/A$ is a root.
\end{theorem}

\begin{proof}
If the three conditions hold, the first two potentially growing lines of
\eqref{eq:master-asymptotic} vanish and all remaining terms tend to zero; hence
$\log\cF_{P,n}\to0$.

Conversely, suppose $\log\cF_{P,n}\to0$.  Since $M>\lambda>1$, a nonzero
$nP(M)M^n$ term cannot be canceled by any later term, so $P(M)=0$.  The remaining bulk term
is $\log A\,MP'(M)M^n$; it forces $P'(M)=0$.  The only remaining growing mode is
$-P(\lambda)\lambda^n/2$, so $P(\lambda)=0$.
\end{proof}

\begin{observation}
The theorem is a classification, not merely a construction.  Every finite-shift factorial
ratio converging to one must use exactly the same three annihilations.  A different-looking
fixed product cannot evade them.
\end{observation}

\subsection{The primitive minimal annihilator}

Let
\begin{equation}
 g:=\gcd(A,M^2),
 \qquad a_0:=\frac Ag,
 \qquad b_0:=\frac{M^2}{g}.
 \label{eq:g-a0-b0}
\end{equation}
Then $\gcd(a_0,b_0)=1$ and $\lambda=b_0/a_0$.  Define
\begin{equation}
 P_*(T):=(T-M)^2(a_0T-b_0).
 \label{eq:Pstar}
\end{equation}
In coefficient order from low to high,
\begin{equation}
 P_*(T)=
 -\frac{M^4}{g}
 +\frac{M^2(A+2M)}{g}T
 -\frac{M^2+2AM}{g}T^2
 +\frac AgT^3.
 \label{eq:Pstar-coefficients}
\end{equation}

\begin{proposition}[Minimality and divisibility]
\label{prop:Pstar-minimal}
The polynomial $P_*$ is primitive.  A polynomial $P\in\Z[T]$ satisfies
\eqref{eq:three-annihilation-conditions} if and only if
\begin{equation}
 P_*(T)\mid P(T)\quad\text{in }\Z[T].
 \label{eq:Pstar-divides}
\end{equation}
In particular, every nonzero fixed factorial near-unit has degree at least three.
\end{proposition}

\begin{proof}
The factors $(T-M)^2$ and $a_0T-b_0$ are primitive and coprime in $\Q[T]$ because
$M\ne M^2/A$.  The rational-root theorem and Gauss's lemma give the claimed divisibility.
Primitivity follows again from Gauss's lemma.
\end{proof}

\begin{proposition}[Sharp defect of the minimal cocycle]
\label{prop:minimal-defect}
Put
\begin{equation}
 C_*:=\frac{M^2(A-M)}{6gA}
       \left(M-\frac{M^3}{A^2}\right)^2>0.
 \label{eq:Cstar}
\end{equation}
Then
\begin{equation}
 \log\cF_{P_*,n}
 =C_*\eta^n+O(q^n),
 \qquad
 \cF_{P_*,n}-1\sim C_*\eta^n.
 \label{eq:minimal-asymptotic}
\end{equation}
In particular, $\cF_{P_*,n}>1$ for all sufficiently large $n$.
\end{proposition}

\begin{proof}
The three growing terms vanish.  Since $\eta>q>\tau$, the leading remaining term in
\eqref{eq:master-asymptotic} is $-P_*(\eta)\eta^n/6$.  Now
\[
 a_0\eta-b_0
 =\frac Ag\frac{M^3}{A^2}-\frac{M^2}{g}
 =-\frac{M^2(A-M)}{gA},
\]
so $-P_*(\eta)/6=C_*$.  Finally $e^u-1\sim u$ as $u\to0$.
\end{proof}

\subsection{Three sign changes are unavoidable}

\begin{proposition}[Sign-variation obstruction]
\label{prop:sign-variation}
Every nonzero polynomial $P\in\Z[T]$ satisfying
\eqref{eq:three-annihilation-conditions} has at least three sign variations in its sequence
of nonzero coefficients.  In particular it has both positive and negative coefficients.
\end{proposition}

\begin{proof}
The polynomial has the three positive roots $M,M,M^2/A$, counted with multiplicity.
Descartes' rule of signs gives at least three coefficient sign variations.
\end{proof}

This simple fact is arithmetically decisive.  Near-unit cancellation forces reciprocals of
some factorial intervals, and those negative intervals expose primes which positive intervals
must somehow transport across geometric scales.

\section{Exact denominator valuations}
\label{sec:valuations}

\subsection{The floor-sum formula}

For a positive rational $y$, let $\den(y)$ be the positive denominator of $y$ in lowest
terms.  For $t\ge0$ and a prime $p$ define
\begin{equation}
 V_p(t):=v_p(F_t)
 =v_p(A^t!)-v_p((A^t-M^t)!).
 \label{eq:Vp-t}
\end{equation}
Legendre's formula gives
\begin{equation}
 V_p(t)=\sum_{r\ge1}
 \left(
   \left\lfloor\frac{A^t}{p^r}\right\rfloor
  -\left\lfloor\frac{A^t-M^t}{p^r}\right\rfloor
 \right).
 \label{eq:Vp-floor}
\end{equation}
The inner difference counts multiples of $p^r$ in the terminal interval
\begin{equation}
 I_t:=(A^t-M^t,A^t].
 \label{eq:It}
\end{equation}

\begin{proposition}[Exact prime-power valuation of a cocycle]
\label{prop:cocycle-valuation}
For $P(T)=\sum_{j=0}^dc_jT^j$,
\begin{align}
 v_p(\cF_{P,n})
 &=\sum_{j=0}^dc_jV_p(n+j)\notag\\
 &=\sum_{r\ge1}\sum_{j=0}^dc_j
 \left(
   \left\lfloor\frac{A^{n+j}}{p^r}\right\rfloor
  -\left\lfloor\frac{A^{n+j}-M^{n+j}}{p^r}\right\rfloor
 \right).
 \label{eq:cocycle-valuation}
\end{align}
Consequently
\begin{equation}
 \log\den(\cF_{P,n})
 =\sum_p\max\{0,-v_p(\cF_{P,n})\}\log p.
 \label{eq:den-log}
\end{equation}
\end{proposition}

\begin{proof}
The first identity is additivity of $p$-adic valuation, the second is Legendre's formula, and
the last is the prime factorization of a reduced denominator.
\end{proof}

\begin{remark}[Finite feasibility]
For fixed $M,A,n,d$, all sums in \eqref{eq:cocycle-valuation} are finite and exact.  Thus the
conditions
\begin{equation}
 \sum_jc_jV_p(n+j)\ge0\qquad(p\le A^{n+d})
 \label{eq:integrality-linear-inequalities}
\end{equation}
are a finite system of homogeneous integer linear inequalities in the coefficient vector
$(c_0,\ldots,c_d)$.  The archimedean annihilations are three homogeneous linear equations
after clearing the denominator of $M^2/A$.  This will become the finite search target in
Section~\ref{sec:search-target}.
\end{remark}

\subsection{Why the visible powers of \texorpdfstring{$A$}{A} cancel}

The normalized probability $R_n$ contains the large denominator $(A^n)^{M^n}$.  For the
minimal shift this visible denominator cancels exactly rather than merely asymptotically.
More generally, if $P(M)=P'(M)=0$, then Lemma~\ref{lem:shift-modes} gives
\begin{equation}
 \sum_jc_j(n+j)M^{n+j}=0.
 \label{eq:A-power-cancellation}
\end{equation}
Therefore
\begin{equation}
 \prod_jR_{n+j}^{c_j}=\prod_jF_{n+j}^{c_j}=\cF_{P,n}.
 \label{eq:R-to-F-cancellation}
\end{equation}
The obstruction is not the obvious power of $A$.  It is the primitive prime content of the
terminal factorial intervals themselves.

\section{Prime windows and the denominator theorem}
\label{sec:primewindows}

\subsection{External analytic inputs}

Write
\begin{equation}
 \vartheta_1(X):=\sum_{p\le X}\log p
 \label{eq:chebyshev-theta}
\end{equation}
for the first Chebyshev function.  We use two classical results.

\begin{enumerate}
\item \textbf{Huxley short-interval prime number theorem.}
For every fixed $\gamma>7/12$,
\begin{equation}
 \vartheta_1(X)-\vartheta_1(X-X^\gamma)\sim X^\gamma.
 \label{eq:huxley}
\end{equation}
The usual statement for $\psi$ implies this form because prime powers contribute
$o(X^\gamma)$ when $\gamma>1/2$ \cite{Huxley1972}.

\item \textbf{Brun--Titchmarsh.}
Uniformly for $2\le h<y$,
\begin{equation}
 \pi(y+h)-\pi(y)\le\frac{2h}{\log h}.
 \label{eq:BT}
\end{equation}
Consequently, if $h\ge X^\delta$ and the interval lies below $O(X^C)$ for fixed $C$, its
logarithmic prime mass is at most $(2C/\delta+o(1))h$\cite{MontgomeryVaughan2007}.
\end{enumerate}

The exponent in our short interval is exactly
$\alpha=\log2/\log3=0.6309\ldots>7/12$, so Huxley's theorem applies with room to spare.

\subsection{The fixed-window mechanism}

Let $P(T)=\sum_{j=0}^dc_jT^j$ have at least one negative coefficient, and let
\begin{equation}
 k:=\max\{j:c_j<0\}.
 \label{eq:k-highest-negative}
\end{equation}
At level $n+k$ set
\begin{equation}
 X:=A^{n+k},
 \qquad H:=M^{n+k}=X^\alpha.
 \label{eq:XH}
\end{equation}
Every prime $p\in(X-H,X]$ occurs once in $F_{n+k}$ and in no lower factorial interval once
$n$ is large.  Because $c_k<0$, it contributes to the denominator unless a higher positive
interval $I_{n+j}$ contains a multiple of $p$.

For $j=k+\ell>k$, write
\begin{equation}
 Y_\ell=A^\ell X,
 \qquad W_\ell=M^\ell H.
 \label{eq:YW}
\end{equation}
A prime $p\in(X-H,X]$ divides $F_{n+k+\ell}$ exactly when some integer $m$ satisfies
\begin{equation}
 Y_\ell-W_\ell<mp\le Y_\ell.
 \label{eq:multiple-capture}
\end{equation}
For fixed $m$, this confines $p$ to an interval of length at most
\begin{equation}
 \frac{W_\ell}{m}\le2Hq^\ell
 \label{eq:capture-length}
\end{equation}
for large $n$, since all relevant $m$ are asymptotic to $A^\ell$.  The number of relevant
multipliers is at most
\begin{equation}
 2+2A^\ell\frac HX.
 \label{eq:multiplier-count}
\end{equation}
These two elementary bounds are the geometric core of the prime-transport estimate.

\subsection{A linearly widening no-go theorem}

\begin{theorem}[Prime-window denominator theorem]
\label{thm:prime-window}
Assume that $x\ge12$.  Fix $\varepsilon>0$.  Let
\[
 P_n(T)=\sum_{j=0}^{d_n}c_{n,j}T^j\in\Z[T]
\]
be any sequence of nonzero polynomials, each having a negative coefficient, and let
$k_n=\max\{j:c_{n,j}<0\}$.  Suppose that
\begin{equation}
 d_n-k_n\le
 (\thetaAE-1-\varepsilon)(n+k_n)
 \label{eq:subcritical-span}
\end{equation}
for all sufficiently large $n$.  Then
\begin{equation}
 \log\den(\cF_{P_n,n})
 \ge(-c_{n,k_n})\bigl(\Delta_\varepsilon+o(1)\bigr)M^{n+k_n},
 \label{eq:prime-window-den-lower}
\end{equation}
where one may take
\begin{equation}
 \Delta_\varepsilon
 :=1-
 \frac{8q}
 {(2-\thetaAE+\varepsilon(1-\alpha))(1-q)}
 >0.
 \label{eq:Delta-epsilon}
\end{equation}
In particular, every fixed polynomial with a negative coefficient has
\begin{equation}
 \log\den(\cF_{P,n})\gg_{P,M,A}M^n.
 \label{eq:fixed-den-exponential}
\end{equation}
\end{theorem}

\begin{proof}
Suppress the index $n$ on $k_n,d_n,c_{n,j}$.  Put $N=n+k$, $X=A^N$, $H=M^N$, and
$D=d-k$.  Huxley's theorem gives
\begin{equation}
 \sum_{X-H<p\le X}\log p=(1+o(1))H.
 \label{eq:all-window-mass}
\end{equation}
Call a prime in this interval \emph{covered} if it divides some $F_{N+\ell}$ with
$1\le\ell\le D$ and positive coefficient $c_{n,k+\ell}>0$.

For a fixed $\ell$, conditions \eqref{eq:capture-length} and
\eqref{eq:multiplier-count} cover all such primes by at most
$2+2A^\ell H/X$ intervals, each of length at most $2Hq^\ell$.
The span hypothesis implies
\begin{align}
 \frac{\log(2Hq^\ell)}{\log X}
 &\ge 2-\thetaAE+\varepsilon(1-\alpha)+o(1)
 =:\delta_\varepsilon+o(1)>0.
 \label{eq:capture-exponent}
\end{align}
Brun--Titchmarsh therefore bounds the logarithmic prime mass in each covering interval by
$(4/\delta_\varepsilon+o(1))Hq^\ell$.  Summing over the multipliers and then over $\ell$ gives
\begin{align}
 \text{covered mass}
 &\le
 \left(\frac8{\delta_\varepsilon}+o(1)\right)H
 \sum_{\ell=1}^{D}
 \left(q^\ell+q^NM^\ell\right)\notag\\
 &\le
 \left(\frac{8q}{\delta_\varepsilon(1-q)}+o(1)\right)H.
 \label{eq:covered-mass}
\end{align}
Indeed,
\[
 q^NM^D
 \le\exp(-\varepsilon N\log M),
\]
so the second geometric sum is negligible.

For $x\ge12$, $q\le(2/3)^{12}$, and already at $\varepsilon=0$ the displayed coefficient is
\begin{equation}
 \frac{8(2/3)^{12}}
 {(2-\log_2 3)(1-(2/3)^{12})}
 =0.1497158446\ldots<1.
 \label{eq:numeric-coverage}
\end{equation}
Hence the uncovered primes have logarithmic mass at least
$(\Delta_\varepsilon+o(1))H$.

An uncovered prime occurs once in $F_{N}=F_{n+k}$, occurs in no lower $F_{n+j}$ because
$p>A^{N-1}$ for large $n$, and occurs in no higher factor with positive coefficient by
definition.  There are no higher negative coefficients because $k$ was maximal.  Thus
$v_p(\cF_{P_n,n})\le c_{n,k}<0$.  Summing its contribution to
\eqref{eq:den-log} over the uncovered primes proves
\eqref{eq:prime-window-den-lower}.
\end{proof}

\begin{corollary}[No fixed-shift factorial proof]
\label{cor:no-fixed-factorial}
For every nonzero fixed $P\in\Z[T]$ with
$\cF_{P,n}\to1$, the reduced denominator satisfies
\begin{equation}
 \log\den(\cF_{P,n})\gg M^n.
 \label{eq:no-fixed-factorial}
\end{equation}
In particular, $\cF_{P,n}$ is not an integer for all sufficiently large $n$ and cannot be
closed by the argument ``an integer converging to one is eventually one.''
\end{corollary}

\begin{proof}
Proposition~\ref{prop:sign-variation} supplies a negative coefficient.  A fixed degree
satisfies \eqref{eq:subcritical-span} for every $\varepsilon<\thetaAE-1$ once $n$ is large.
Apply Theorem~\ref{thm:prime-window}.
\end{proof}

\begin{corollary}[Minimal cocycle denominator]
\label{cor:minimal-denominator}
For the canonical polynomial $P_*$,
\begin{equation}
 \cF_{P_*,n}=1+C_*\eta^n+o(\eta^n),
 \qquad
 \log\den(\cF_{P_*,n})\gg M^n.
 \label{eq:minimal-gap-den}
\end{equation}
Thus its archimedean distance from one is merely exponentially small in $n$, whereas its
reduced denominator is exponentially large in $M^n$.
\end{corollary}

\subsection{Two linear phase thresholds}

The proof exposes two distinct universal scales for a higher interval $I_{n+k+\ell}$.

\begin{enumerate}
\item The multiplier $m$ in \eqref{eq:multiple-capture} begins to drift across more than one
integer when
\begin{equation}
 A^\ell\frac HX\asymp1,
 \qquad
 \frac\ell{n+k}\asymp1-\alpha
 =0.3690702464\ldots.
 \label{eq:drift-threshold}
\end{equation}

\item The higher interval itself becomes as long as the prime scale $X$ when
\begin{equation}
 \frac{M^{n+k+\ell}}{A^{n+k}}\asymp1,
 \qquad
 \frac\ell{n+k}\asymp\thetaAE-1
 =0.5849625007\ldots.
 \label{eq:thickness-threshold}
\end{equation}
\end{enumerate}

The multiplier-drift threshold does not suffice: the covering argument survives throughout
the intermediate regime and fails only at the thickness threshold.  Consequently a
factorial-cocycle rescue must be genuinely linearly nonlocal.

\begin{resultbox}
\textbf{Exact escape requirement.}
Any integer-near-one factorial-ratio proof based on shift annihilation must use, beyond its
highest negative factor, positive factors at distance at least
\[
 (\log_2 3-1-o(1))n.
\]
Fixed order, logarithmic order, $n^\gamma$ order with $\gamma<1$, and every linear order below
that constant are excluded.
\end{resultbox}

\section{Binomial cocycles: the same obstruction with two extra roots}
\label{sec:binomial}

The integral sequence
\begin{equation}
 B_n:=\binom{A^n}{M^n}=\frac{F_n}{M^n!}
 \label{eq:Bn}
\end{equation}
is an obvious attempt to remove some denominator anxiety at the outset.  Signed products of
$B_n$ still have rational denominators, but every individual factor is integral.  The
archimedean classification changes in a precise way.

\begin{proposition}[Binomial asymptotic]
\label{prop:binomial-asymptotic}
As $n\to\infty$,
\begin{align}
 \log B_n
 &=nM^n\log\frac AM+M^n
   -\frac12n\log M-\frac12\log(2\pi)
   -\frac12\lambda^n+\frac12q^n
   -\frac16\eta^n
   +O(\tau^n+M^{-n}).
 \label{eq:binomial-asymptotic}
\end{align}
\end{proposition}

\begin{proof}
Subtract Stirling's expansion
\[
 \log(M^n!)=nM^n\log M-M^n+\frac12n\log M+\frac12\log(2\pi)+O(M^{-n})
\]
from Proposition~\ref{prop:Fn-expansion}.
\end{proof}

For $P(T)=\sum c_jT^j$, put
\begin{equation}
 \cB_{P,n}:=\prod_jB_{n+j}^{c_j}.
 \label{eq:binomial-cocycle}
\end{equation}

\begin{theorem}[Complete binomial near-unit classification]
\label{thm:binomial-classification}
For fixed $P\in\Z[T]$,
\begin{equation}
 \cB_{P,n}\longrightarrow1
 \label{eq:binomial-nearunit}
\end{equation}
if and only if
\begin{equation}
 P(M)=P'(M)=P(1)=P'(1)=P(M^2/A)=0.
 \label{eq:binomial-roots}
\end{equation}
The primitive minimal annihilator is
\begin{equation}
 P_{\mathrm{bin}}(T)=(T-1)^2P_*(T).
 \label{eq:Pbin}
\end{equation}
\end{theorem}

\begin{proof}
Apply Lemma~\ref{lem:shift-modes} to
\eqref{eq:binomial-asymptotic}.  The $nM^n$ and $M^n$ terms force a double root at $M$.
The $\lambda^n$ term forces the root $M^2/A$.  Finally,
\[
 P(\EF)n=nP(1)+P'(1),
\]
so the logarithmic Stirling term forces a double root at $1$.  These conditions also suffice.
The divisibility statement follows from Gauss's lemma.
\end{proof}

\begin{theorem}[Binomial prime-window obstruction]
\label{thm:binomial-window}
Under the hypotheses of Theorem~\ref{thm:prime-window}, the same denominator lower bound
holds with $\cF_{P_n,n}$ replaced by $\cB_{P_n,n}$.
\end{theorem}

\begin{proof}
For an uncovered prime $p\asymp A^{n+k_n}$, the span condition gives
\[
 M^{n+d_n}/p
 \le q^{n+k_n}M^{d_n-k_n}
 \le e^{-\varepsilon(n+k_n)\log M}\to0.
\]
Thus $p$ exceeds every lower factorial $M^{n+j}!$ occurring in the binomial denominators.
Its valuations in the $B_{n+j}$ therefore coincide with its valuations in the terminal
products $F_{n+j}$.  The proof of Theorem~\ref{thm:prime-window} applies unchanged.
\end{proof}

The binomial version removes the visible nonintegrality of each building block but does not
remove the signed prime transport forced by archimedean cancellation.  It also costs two
extra positive roots and therefore at least two additional coefficient sign variations in
typical minimal constructions.

\section{Higher collision annihilation and why it is not a free rescue}
\label{sec:higher-collision}

The minimal polynomial cancels the pair-collision mode $\lambda^n$ and leaves the
triple-collision mode $\eta^n$.  It is natural to cancel further modes.

\subsection{The full mode lattice}

From the absolutely convergent identity
\begin{equation}
 \log R_n=-\sum_{r\ge1}\frac1{rA^{rn}}
              \sum_{j=0}^{M^n-1}j^r,
 \label{eq:full-log-series}
\end{equation}
Faulhaber's formula shows that every fixed collision order $r$ contributes a finite linear
combination of geometric modes
\begin{equation}
 \rho_{r,s}^n,
 \qquad
 \rho_{r,s}:=\frac{M^s}{A^r},
 \qquad 1\le s\le r+1.
 \label{eq:mode-lattice}
\end{equation}
The largest mode at order $r$ is
\begin{equation}
 \rho_{r,r+1}=\frac{M^{r+1}}{A^r}
 =\left(2\left(\frac23\right)^r\right)^x.
 \label{eq:leading-r-mode}
\end{equation}
Only the case $r=1$ is greater than one.  The descending subunit modes begin
\begin{equation}
 \frac{M^3}{A^2}=\left(\frac89\right)^x,
 \quad
 \frac MA=\left(\frac23\right)^x,
 \quad
 \frac{M^4}{A^3}=\left(\frac{16}{27}\right)^x,
 \quad
 \frac{M^2}{A^2}=\left(\frac49\right)^x,
 \quad\ldots.
 \label{eq:mode-order}
\end{equation}

\subsection{Finite high-order annihilators}

For a finite set $S$ of rational modes, one may form the primitive integer polynomial
\begin{equation}
 P_S(T)=(T-M)^2\prod_{\rho=u/v\in S}(vT-u),
 \label{eq:PS}
\end{equation}
after removing duplicate factors and common content.  If $M^2/A\in S$, the associated
factorial cocycle is a near-unit; each additional root pushes its leading asymptotic to the
next unannihilated mode.

This does not address the denominator theorem.  Every nonzero $P_S$ still has mixed signs,
and every fixed set $S$ has fixed degree.  Corollary~\ref{cor:no-fixed-factorial} therefore
forces denominator logarithm $\gg M^n$ regardless of how many collision modes have been
canceled.

\begin{warning}[Cumulant cancellation versus arithmetic cancellation]
\label{warn:cumulant-arithmetic}
Adding roots can make $|\cF_{P,n}-1|$ spectacularly small while making the coefficient vector
and denominator more complicated.  It does not turn the rational near-unit into an integer.
The archimedean and prime-window problems are independent constraints and must be solved
simultaneously.
\end{warning}

\subsection{A useful critical coincidence}

If the triple mode $\eta$ is also canceled, the next mode is $q=M/A$ unless it too is
annihilated.  Its decay rate is
\begin{equation}
 -\log q=x\log\frac32.
 \label{eq:q-decay}
\end{equation}
The coefficient-growth required to move factorial mass across the thickness threshold also
naturally involves $\log(A/M)=x\log(3/2)$.  Thus the first higher-order refinement lands on an
exact critical boundary rather than creating abundant slack.  This coincidence is a warning
against naive ``cancel one more term'' optimism, but it also suggests that a sharp extremal
construction might exist exactly at the boundary.

\section{The remaining linear-width problem}
\label{sec:search-target}

Theorem~\ref{thm:prime-window} does more than reject a construction.  It identifies the
first scale at which a construction could possibly work.

\subsection{Exact finite formulation}

For fixed $M,A,n,d$, seek a nonzero vector
\begin{equation}
 c=(c_0,\ldots,c_d)\in\Z^{d+1}
 \label{eq:c-vector}
\end{equation}
subject to the three archimedean equations
\begin{align}
 \sum_{j=0}^dc_jM^j&=0,
 \label{eq:lin-M}\\
 \sum_{j=0}^djc_jM^j&=0,
 \label{eq:lin-derivative}\\
 \sum_{j=0}^dc_jM^{2j}A^{d-j}&=0.
 \label{eq:lin-lambda}
\end{align}
The last equation is $A^dP(M^2/A)=0$.  For integrality impose, for every prime
$p\le A^{n+d}$,
\begin{equation}
 \sum_{j=0}^dc_jV_p(n+j)\ge0.
 \label{eq:prime-ILP}
\end{equation}
A solution makes $\cF_{P,n}$ an integer and cancels all growing archimedean modes.
For a family $c=c^{(n)}$, one must additionally control the shifted remainder, for example
through
\begin{equation}
 \sum_{j=0}^{d_n}|c_{n,j}|\eta^j=o(\eta^{-n})
 \label{eq:height-condition-crude}
\end{equation}
or, more sharply, through direct evaluation of $P_n$ on the ordered mode lattice.

\begin{target}[Linear-width factorial feasibility]
\label{target:linear-width}
Decide whether there is a sequence $P_n\in\Z[T]$ such that:
\begin{enumerate}
\item $P_n(M)=P_n'(M)=P_n(M^2/A)=0$;
\item $v_p(\cF_{P_n,n})\ge0$ for every prime $p$;
\item $\cF_{P_n,n}\to1$;
\item $\deg P_n-k_n\ge(\log_2 3-1-o(1))(n+k_n)$, where $k_n$ is the highest negative index.
\end{enumerate}
A positive answer, together with an eventual strict inequality
$\cF_{P_n,n}\ne1$ (for example from a sign-controlled first surviving archimedean mode),
would prove the conjecture: a sequence of positive integers converging to $1$ is eventually
identically $1$.  A negative answer, proved by a dual valuation certificate, would close the
entire factorial-cocycle family.
\end{target}

\subsection{Valuation cone and dual certificates}

Let
\begin{equation}
 \mathbf v_p^{(n,d)}=(V_p(n),V_p(n+1),\ldots,V_p(n+d))\in\Nzero^{d+1}.
 \label{eq:valuation-vector}
\end{equation}
The integrality region is the rational polyhedral cone
\begin{equation}
 \mathcal C_{n,d}:=
 \{c\in\R^{d+1}:\langle c,\mathbf v_p^{(n,d)}\rangle\ge0
 \text{ for all }p\le A^{n+d}\}.
 \label{eq:valuation-cone}
\end{equation}
The archimedean constraints define a codimension-three subspace
$\mathcal L_{M,A,d}$.  The finite question is whether
\begin{equation}
 \mathcal C_{n,d}\cap\mathcal L_{M,A,d}
 \label{eq:cone-intersection}
\end{equation}
contains a nonzero integer point satisfying a height bound.

The proof of Theorem~\ref{thm:prime-window} is already a dual certificate below the critical
width: a positive weighted combination of the prime inequalities detects the highest
negative coefficient.  At and above the critical width, the right search is therefore not
blind enumeration of $c$ but computation of sparse dual certificates supported on prime
windows and their transported multiples.

\subsection{Three increasingly ambitious search layers}

\begin{enumerate}
\item \textbf{LP layer.}  Ignore integrality of $c$, normalize one linear functional, and test
whether the real cone intersection is empty.  An infeasible LP returns a rational Farkas
certificate which can be verified exactly.

\item \textbf{ILP layer.}  If the LP is feasible, bound $\|c\|_1$ or the weighted height
$\sum|c_j|\eta^j$ and search for integer points.  Root divisibility by $P_*$ can reduce the
variables by writing $P_n=P_*S_n$.

\item \textbf{Prime-power layer.}  Initial experiments may impose only primes
$p>A^{n+d}/2$ and then descend through dyadic prime ranges.  Once primes are exhausted, add
$p^r$ constraints from \eqref{eq:cocycle-valuation}; these are essential for actual
integrality but not for the first window obstruction.
\end{enumerate}

\subsection{Primitive-output sensitivity}

A successful proof must eventually distinguish hypothetical primitive outputs from the
legitimate controls $(M,A)=(2^m,3^m)$.  The factorial asymptotic alone does not do so: it uses
only the fixed ratio $\log M/\log A$.  The inherited finite exclusion supplies one useful
asymmetry, namely $q\le(2/3)^{12}$, but integer controls with $m\ge12$ also satisfy it.

The most promising primitive-sensitive refinement is to add local constraints at a
structural prime $r\mid M$ with $r\ne2$, or at a prime $s\mid A$ with $s\ne3$.  Such a prime
must exist for a nonintegral solution: if $M$ were a pure power of $2$, then
$2^x=M=2^m$ would force $x=m$; similarly for $A$ and powers of $3$.  For $r\mid M$ and
$r\nmid A$, Legendre's formula gives an exact low-digit cancellation in
$\binom{A^n}{M^n}$ for all powers $r^a\mid M^n$.  This creates a long carry-free prefix in
base $r$ and may be combined with the high-prime transport cone.  I did not complete that
hybrid argument, but it is the first remaining direction that genuinely uses nonintegrality
rather than merely the ratio $\log2/\log3$.

\begin{target}[Hybrid high-prime/structural-prime certificate]
\label{target:hybrid}
Construct a dual certificate combining:
\begin{enumerate}
\item high primes in the terminal interval $(A^N-M^N,A^N]$, which constrain coefficient
      signs and shift span;
\item a structural prime dividing the odd part of $M$ or the non-$3$ part of $A$, which
      constrains the same coefficients through long base-$p$ carry-free prefixes.
\end{enumerate}
The desired contradiction is that the high-prime cone requires prime mass to move forward by
at least $(\log_2 3-1)N$ shifts, while the structural-prime valuation forces the corresponding
coefficient mass to remain locally balanced.
\end{target}

\section{Exact computational audit}
\label{sec:computation}

The analytic argument does not depend on floating-point experimentation.  Computation is
nevertheless useful for checking signs, transients, valuation formulas, and denominator
scale.  The legitimate control $(M,A)=(2,3)$ is especially useful: it satisfies the same
mesoscopic identities but cannot lead to a contradiction because $x=1$ is an allowed
integer exponent.

For this control,
\begin{equation}
 P_*(T)=-16+28T-16T^2+3T^3.
 \label{eq:Pstar-23}
\end{equation}
The table computes $\log\cF_{P_*,n}$ by high-precision log-gamma evaluation and computes the
reduced denominator exactly, prime by prime, using Legendre valuations.  The denominator
column is $\log\den(\cF_{P_*,n})$.

\begin{table}[ht]
\centering
\small
\begin{tabular}{r@{\qquad}r@{\qquad}r@{\qquad}r}
\toprule
$n$ & $\log\cF_{P_*,n}$ & $\log|\cF_{P_*,n}-1|$ & $\log\den(\cF_{P_*,n})$\\
\midrule
2  & $-0.808212$ & $-0.586617$ & $386.87$\\
3  & $-0.408119$ & $-1.095066$ & $914.74$\\
4  & $-0.197859$ & $-1.718805$ & $2299.56$\\
5  & $-0.082391$ & $-2.537378$ & $5332.50$\\
6  & $-0.018582$ & $-3.994918$ & $11795.74$\\
7  & $ 0.019229$ & $-3.941309$ & $25135.36$\\
8  & $ 0.040066$ & $-3.177793$ & $52856.16$\\
9  & $ 0.048687$ & $-2.997675$ & $109568.88$\\
10 & $ 0.048779$ & $-2.995638$ & $224549.84$\\
\bottomrule
\end{tabular}
\caption{Minimal factorial cocycle for the integer control $(M,A)=(2,3)$.  The eventual
positive sign predicted by Proposition~\ref{prop:minimal-defect} appears at $n=7$.  The
asymptotic decay is still in a long transient, while the exact reduced denominator is already
enormous.}
\label{tab:control-audit}
\end{table}

The sign transition is not a defect in the proof: the leading positive $\eta^n$ term must
first overtake the negative $q^n$ and smaller modes.  More importantly, the denominator is
not an artifact of retaining the visible powers $(A^n)^{M^n}$; those powers cancel exactly
by \eqref{eq:A-power-cancellation}.  The table measures only prime factors in the terminal
factorial intervals.

\subsection{Reproducible valuation code}

The following compact script reproduces the denominator column.  It avoids constructing the
astronomically large numerator and denominator; all arithmetic defining the reduced
denominator is exact.

\begin{lstlisting}[style=pythonstyle,caption={Exact valuation audit for a factorial cocycle.}]
from math import log
from mpmath import mp
from sympy import primerange


def vp_factorial(n: int, p: int) -> int:
    value = 0
    while n:
        n //= p
        value += n
    return value


def falling_factorial_valuation(t: int, M: int, A: int, p: int) -> int:
    top = A**t
    width = M**t
    return vp_factorial(top, p) - vp_factorial(top - width, p)


def p_star_coefficients(M: int, A: int) -> list[int]:
    from math import gcd
    g = gcd(A, M*M)
    # Coefficients from T^0 through T^3.
    return [
        -(M**4)//g,
        (M*M*(A + 2*M))//g,
        -(M*M + 2*A*M)//g,
        A//g,
    ]


def log_falling_factorial(t: int, M: int, A: int):
    top = A**t
    width = M**t
    return mp.loggamma(top + 1) - mp.loggamma(top - width + 1)


def audit(M: int, A: int, n: int):
    coefficients = p_star_coefficients(M, A)
    degree = len(coefficients) - 1

    mp.dps = 80
    log_q = mp.fsum(
        coefficients[j] * log_falling_factorial(n + j, M, A)
        for j in range(degree + 1)
    )
    log_gap = mp.log(abs(mp.expm1(log_q)))

    log_denominator = mp.mpf("0")
    for p in primerange(2, A**(n + degree) + 1):
        valuation = sum(
            coefficients[j]
            * falling_factorial_valuation(n + j, M, A, p)
            for j in range(degree + 1)
        )
        if valuation < 0:
            log_denominator += (-valuation) * mp.log(p)

    return log_q, log_gap, log_denominator


for n in range(2, 11):
    print(n, audit(M=2, A=3, n=n))
\end{lstlisting}

\subsection{Finite certificate mode}

For proof-oriented computation, numerical logarithms should be removed.  The exact kernel is:
\begin{enumerate}
\item generate each prime $p\le A^{n+d}$ together with a primality certificate;
\item compute $V_p(n+j)$ by integer division;
\item evaluate the dot product $\sum_jc_jV_p(n+j)$;
\item record negative valuations as a factorized denominator certificate;
\item bound $\log\cF_{P,n}$ using the rational inequalities in
      Lemma~\ref{lem:log-remainder}, rather than log-gamma.
\end{enumerate}
Every output can then be checked by a small independent verifier or imported into Lean.

\section{Lean formalization plan}
\label{sec:lean}

The new argument separates cleanly into an elementary formal core and an analytic-number-
theory interface.

\subsection{Elementary core}

A practical file decomposition is:

\begin{description}[leftmargin=3.4cm,style=nextline]
\item[\texttt{BirthdayProduct.lean}]
Define falling factorials, $F_n$, $R_n$, and prove positivity and the product formula.
Formalize Lemma~\ref{lem:log-remainder} and the two power-sum identities.

\item[\texttt{MesoscopicBounds.lean}]
Prove $1/2<\log2/\log3<2/3$ from $4>3$ and $8<9$, then derive
$M^2/A>1$, $M^3/A^2<1$, and the ordering of $q,\eta,\tau$ under
$M=2^x$, $A=3^x$.

\item[\texttt{ShiftPolynomial.lean}]
For a finitely supported integer coefficient function, define $P(\EF)$ and prove
\eqref{eq:shift-exp}--\eqref{eq:shift-nexp}.  Connect the implementation to
\texttt{Polynomial.eval} and formal derivatives.

\item[\texttt{FactorialCocycle.lean}]
Define $\cF_{P,n}$ in $\Q_{>0}$ and prove the near-unit sufficiency bounds.
For necessity, use explicit lower bounds separating $nM^n$, $M^n$, and $\lambda^n$.

\item[\texttt{MinimalAnnihilator.lean}]
Prove divisibility by $(T-M)^2(a_0T-b_0)$ using polynomial division and Gauss content.
The asymptotic sign can be formalized as an eventual inequality without developing a general
asymptotics framework.

\item[\texttt{FactorialValuation.lean}]
Prove Legendre's floor formula, \eqref{eq:cocycle-valuation}, and the reduced-denominator
identity.  This part should reuse Mathlib's factorial and padic valuation infrastructure where
available, with a local terminal-interval count lemma to avoid normalization friction.
\end{description}

\subsection{Prime-window interface}

Formalizing Huxley's theorem from first principles is far beyond the sensible scope of this
project.  Isolate it behind one statement:
\begin{equation}
 \forall\gamma>7/12,\quad
 \sum_{X-X^\gamma<p\le X}\log p=(1+o(1))X^\gamma.
 \label{eq:lean-huxley-interface}
\end{equation}
The Brun--Titchmarsh input should similarly be isolated.  Everything from those two
interfaces to Theorem~\ref{thm:prime-window} is elementary real inequalities, interval
covering, geometric sums, and valuation bookkeeping.

For a fully kernel-checked finite result, replace Huxley with an explicit list of primes in
the relevant window and replace Brun--Titchmarsh with direct enumeration.  This does not
prove the asymptotic theorem, but it validates the combinatorial transport layer and the ILP
encoding.

\subsection{Suggested theorem order}

The order that minimizes dependency churn is:
\begin{enumerate}
\item exact log inequalities and power sums;
\item shift identities;
\item explicit upper and lower bounds for $\log\cF_{P,n}$;
\item root necessity and minimal polynomial;
\item valuation floor sums;
\item fixed-degree prime window using an abstract short-interval-prime hypothesis;
\item linearly widening covering theorem;
\item finite cone/dual-certificate checker.
\end{enumerate}

The exact polynomial identities and valuation formulas are suitable early formalization
targets even before the external analytic interface is settled.

\section{What this changes in the global proof search}
\label{sec:global-assessment}

\subsection{Progress obtained}

\begin{enumerate}
\item \textbf{A genuinely new invariant.}
The terminal factorial interval converts the fixed exponent
$\log2/\log3$ into a one-growing-mode collision expansion.

\item \textbf{A complete local theory.}
Every fixed-shift near-unit is classified by three polynomial roots, and the minimal one is
explicit.

\item \textbf{A quantified arithmetic obstruction.}
The denominator is not merely ``large'' in experiments.  Its logarithm is bounded below by
a positive multiple of $M^n$ using uncaptured short-interval primes.

\item \textbf{A sharp scale for any rescue.}
A successful factorial construction must widen linearly past
$(\log_2 3-1)n$.  This excludes an enormous design space at once.

\item \textbf{A finite next problem.}
Above the threshold, the remaining question is an exact cone intersection with three linear
root equations and finitely many valuation inequalities.
\end{enumerate}

\subsection{What has not been proved}

No contradiction has yet been extracted from the existence of $M,A$.  A rational number may
be extremely close to one when its denominator is enormous, and
Corollary~\ref{cor:minimal-denominator} shows that the minimal cocycle has far more denominator
than the elementary rational-separation bound needs.  The result is therefore not a proof in
disguise.

The unresolved bridge is one of the following:
\begin{enumerate}
\item find a linearly widening cocycle that is integral and still tends to one;
\item prove, by a global dual certificate, that no such cocycle can exist;
\item combine high-prime transport with a primitive structural prime to force incompatible
      coefficient balances;
\item discover a different integer-valued transform of the same birthday expansion whose
      local prime windows cancel automatically.
\end{enumerate}

\subsection{Why the result is still strategically useful}

The supplied report contains many methods whose bottleneck is ``denominator control'' in a
broad sense.  Here the denominator has been localized: it comes from primes in explicit
terminal intervals, and their possible cancellation is a geometric covering problem.  The
phrase ``need better denominator control'' has been replaced by a numerical boundary, an
exact valuation matrix, and a dual-certificate program.  That is a materially narrower
research problem.

\section{Prioritized next experiments}
\label{sec:next}

\begin{enumerate}
\item \textbf{LP feasibility at the threshold.}
For the controls $(M,A)=(2,3)$ and $(4,9)$, test degrees
$d=\lceil(\log_2 3-1)n\rceil+r$ for small offsets $r$.  Use only high-prime valuation rows at
first.  Record Farkas certificates, not just solver status.

\item \textbf{Root-factor parameterization.}
Write $P=P_*S$ and formulate valuations directly in the coefficients of $S$.  Compare sparse,
cyclotomic, and alternating $S$ families.  Add roots at $\eta$ and $q$ only when the measured
coefficient height warrants them.

\item \textbf{Structural-prime rows.}
Introduce symbolic rows corresponding to a prime $r\mid M$, $r\ne2$, in the cases
$r\nmid A$ and $r\mid A$.  Search for a dual combination with high-prime rows before doing
full Kummer digit analysis.

\item \textbf{Continuous transport relaxation.}
Scale shift index by $n$ and terminal depth by $M^n$.  Derive the limiting measure-transport
problem for positive and negative coefficient mass.  Theorem~\ref{thm:prime-window} supplies
the subcritical forbidden region; the goal is a continuum dual functional at and above the
critical width.

\item \textbf{Lean core.}
Formalize Proposition~\ref{prop:Fn-expansion}, Theorem~\ref{thm:near-unit-classification},
and Proposition~\ref{prop:cocycle-valuation}.  These are exact, self-contained, and likely to
be reused even if the final proof leaves the factorial route.
\end{enumerate}

\section{Conclusion}
\label{sec:conclusion}

The new factorial route begins with a favorable accident that is specific to the consecutive
bases two and three:
\[
 \frac12<\frac{\log2}{\log3}<\frac23.
\]
It turns the terminal falling factorial $(A^n)_{M^n}$ into an asymptotic object with exactly
one divergent collision mode.  Finite shifts can annihilate that mode and the bulk term
perfectly, producing a canonical rational number converging exponentially to one.  The
classification theorem shows that this is not one example among many; it is the universal
shape of every fixed-shift factorial near-unit.

The arithmetic audit is equally rigid.  Negative coefficients are forced, and their terminal
intervals contain short-interval primes that higher positive intervals cannot absorb below a
linear transport width.  With $x\ge12$, Brun--Titchmarsh covering leaves at least about
$85\%$ of the logarithmic prime mass uncaptured even at the limiting constant used in the
proof.  Huxley's theorem then forces an exponentially large reduced denominator.  The
integer-near-one shortcut is therefore impossible for all fixed and subcritical-width
factorial and binomial cocycles.

For the integer-near-one factorial architecture, the surviving target is sharp:
\[
 \text{shift width}\ \ge
 (\log_2 3-1-o(1))n.
\]
Above that boundary, archimedean annihilation and prime-power integrality become one finite
integer linear problem.  Either a feasible family yields a proof, or a dual certificate can
close a large new class of approaches.  The immediate next move should be to compute and
formalize that boundary problem rather than to invent another unquantified denominator
heuristic.

\appendix

\section{The canonical cocycle written out}
\label{app:canonical}

With $g=\gcd(A,M^2)$, the minimal polynomial coefficients from
\eqref{eq:Pstar-coefficients} give
\begin{align}
 \cF_{P_*,n}
 &=F_n^{-M^4/g}
   F_{n+1}^{M^2(A+2M)/g}
   F_{n+2}^{-(M^2+2AM)/g}
   F_{n+3}^{A/g}.
 \label{eq:canonical-expanded}
\end{align}
Writing out $F_t=(A^t)!/(A^t-M^t)!$ produces a pure ratio of eight ordinary factorials.
There is no hidden nonintegral exponent.

The cancellation of the powers used in $R_t$ can also be checked directly.  The exponent of
$A$ in
$\prod_{j=0}^3(A^{n+j})^{c_jM^{n+j}}$ is
\begin{equation}
 \sum_{j=0}^3c_j(n+j)M^{n+j}
 =M^n\bigl(nP_*(M)+MP_*'(M)\bigr)=0.
 \label{eq:canonical-A-cancel}
\end{equation}
The remaining distance from one is therefore entirely due to the finite-product corrections
$1-j/A^t$.

For the unreduced but visually simpler annihilator
\begin{equation}
 \widetilde P(T)=(T-M)^2(AT-M^2)=gP_*(T),
 \label{eq:Ptilde}
\end{equation}
the first cancellation may be displayed as
\begin{equation}
 \frac{R_{n+1}^{A}}{R_n^{M^2}},
 \label{eq:first-cancel}
\end{equation}
which kills the pair-collision mode because
$A(M^2/A)^{n+1}=M^2(M^2/A)^n$.  Applying $(\EF-M)^2$ then kills both
$nM^n$ and $M^n$ bulk modes.  This factorization is useful for intuition, while $P_*$ is the
right primitive arithmetic normalization.

\section{An elementary fixed-degree prime slice}
\label{app:fixed-slice}

The full covering theorem is needed for linearly growing degrees.  For a fixed polynomial one
can see the obstruction in one contiguous interval.

Let $k$ be the highest negative coefficient and define
\begin{equation}
 \rho:=\max\left(
 \left\{\left(\frac MA\right)^{j-k}:j>k,\ c_j>0\right\}\cup\{0\}
 \right).
 \label{eq:rho-fixed}
\end{equation}
Then $0\le\rho\le M/A<1$.  Put $X=A^{n+k}$ and $H=M^{n+k}$.  For fixed small
$\epsilon>0$, consider primes
\begin{equation}
 X-(1-\epsilon)H<p<X-(\rho+\epsilon)H.
 \label{eq:fixed-prime-slice}
\end{equation}
Write $p=X-t$.  The only multiple of $p$ that can lie close to the top of the higher interval
$I_{n+j}$ is $A^{j-k}p$, because the interval width is $o(p)$ and
$A^{n+j}/p=A^{j-k}+o(1)$.  That multiple lies in $I_{n+j}$ only if
\begin{equation}
 A^{j-k}t\le M^{n+j},
 \quad\text{equivalently}\quad
 t\le\left(\frac MA\right)^{j-k}H\le\rho H.
 \label{eq:fixed-capture-condition}
\end{equation}
The primes in \eqref{eq:fixed-prime-slice} therefore occur in the negative factor at $k$ and
in no positive factor.  Huxley's theorem gives their logarithmic mass
\begin{equation}
 (1-\rho-2\epsilon+o(1))H.
 \label{eq:fixed-slice-mass}
\end{equation}
Letting $\epsilon\downarrow0$ yields the sharper fixed-degree estimate
\begin{equation}
 \log\den(\cF_{P,n})
 \ge(-c_k)(1-\rho+o(1))M^{n+k}.
 \label{eq:fixed-sharp-den}
\end{equation}
This proof does not use $x\ge12$; that hypothesis is needed only to control the union of many
moving capture intervals in the linearly widening theorem.

\section{A structural-prime carry-free prefix}
\label{app:structural-prime}

The following exact lemma records the primitive-sensitive observation used in
Target~\ref{target:hybrid}.

\begin{lemma}[Low prime powers contribute no binomial carry]
\label{lem:carry-free-prefix}
Let $r$ be a prime with $r\mid M$ and $r\nmid A$, and put $e=v_r(M)$.  For every
$1\le a\le en$,
\begin{equation}
 \left\lfloor\frac{A^n}{r^a}\right\rfloor
 -\left\lfloor\frac{M^n}{r^a}\right\rfloor
 -\left\lfloor\frac{A^n-M^n}{r^a}\right\rfloor=0.
 \label{eq:carry-free-floor}
\end{equation}
Consequently
\begin{equation}
 v_r\binom{A^n}{M^n}
 \le
 \max\left\{0,
 \left\lfloor n\log_r A\right\rfloor-en+1
 \right\}.
 \label{eq:binomial-r-bound}
\end{equation}
\end{lemma}

\begin{proof}
For $a\le en$, the integer $M^n/r^a$ is integral.  Hence
\[
 \left\lfloor\frac{A^n-M^n}{r^a}\right\rfloor
 =\left\lfloor\frac{A^n}{r^a}-\frac{M^n}{r^a}\right\rfloor
 =\left\lfloor\frac{A^n}{r^a}\right\rfloor-\frac{M^n}{r^a},
\]
which proves \eqref{eq:carry-free-floor}.  For $a>en$, each summand in Legendre's formula for
the binomial valuation is either zero or one.  No term remains once $r^a>A^n$, so the number
of possible nonzero terms gives \eqref{eq:binomial-r-bound}.
\end{proof}

In Kummer language, the lowest $en$ base-$r$ digit positions create no carry when adding
$M^n$ and $A^n-M^n$.  The high-prime obstruction is of size $M^n$, while this structural
valuation is only $O(n)$.  That disparity may be useful in a dual argument, but a mechanism
linking the two coefficient weightings is still missing.

\section{Proof-audit checklist}
\label{app:audit}

The following points were checked explicitly because they are common failure modes in this
problem.

\begin{enumerate}
\item \textbf{No four-exponentials assumption.}
All identities use only $M=2^x$, $A=3^x$, elementary inequalities, factorials, and classical
prime-distribution theorems.

\item \textbf{Integer controls survive.}
The construction applies to $(M,A)=(2^m,3^m)$ as well.  It does not claim those pairs are
impossible; it diagnoses why the rational near-unit has enough denominator to coexist with
them.

\item \textbf{No false rational separation.}
The bound $|a/b-1|\ge1/b$ requires an upper bound on $b$ to be useful.  The report proves a
large lower bound instead and explicitly treats it as an obstruction, not a contradiction.

\item \textbf{The shift root at $M$ is double.}
A simple root kills $nM^n$ but leaves the $M^n$ derivative term.  This is why
$P(M)=P'(M)=0$ is necessary.

\item \textbf{The growing collision root is rational.}
$M^2/A$ need not be integral.  The primitive factor is
$(A/g)T-M^2/g$, not necessarily $AT-M^2$.

\item \textbf{Short-interval exponent has margin.}
$\alpha=0.630929\ldots>7/12=0.583333\ldots$, and prime powers are negligible because
$\alpha>1/2$.

\item \textbf{The linearly widening theorem uses the inherited finite exclusion.}
The numerical coverage constant is below one because $q\le(2/3)^{12}$.  The sharper fixed-
degree slice does not require this exclusion.

\item \textbf{Positive higher factors are overcounted.}
The covering proof marks a prime as covered if it divides any positive factor, regardless of
whether the positive coefficient is large enough to cancel $c_k$.  This only weakens the
lower bound and is therefore safe.

\item \textbf{Prime powers are not omitted from exact integrality.}
The finite ILP uses all $p^r$ through Legendre's formula.  The prime-window theorem needs only
first powers because those already certify negative valuation.

\item \textbf{No claim of a completed proof.}
The remaining linear-width feasibility problem and the hybrid structural-prime bridge are
stated as open targets.
\end{enumerate}

\begin{thebibliography}{99}

\bibitem{AlaogluErdos1944}
L. Alaoglu and P. Erd\H{o}s,
\newblock On highly composite and similar numbers,
\newblock \emph{Transactions of the American Mathematical Society} 56 (1944), 448--469.

\bibitem{UnifiedReport2026}
V. Reshetnikov et al.,
\newblock \emph{Unified progress report on the Alaoglu--Erd\H{o}s conjecture},
\newblock supplied manuscript, August 22, 2026.

\bibitem{Huxley1972}
M. N. Huxley,
\newblock On the difference between consecutive primes,
\newblock \emph{Inventiones Mathematicae} 15 (1972), 164--170.

\bibitem{MontgomeryVaughan2007}
H. L. Montgomery and R. C. Vaughan,
\newblock \emph{Multiplicative Number Theory I: Classical Theory},
\newblock Cambridge Studies in Advanced Mathematics 97, Cambridge University Press, 2007.

\end{thebibliography}

\end{document}
